\documentclass{situationist-preprint}

\title{Moderate Moralism Defended}
\msid{AES-08-a}
\msfamily{AES-08}
\msdiscipline{Aesthetics}
\mstype{Principal manuscript}
\mskeywords{ethical criticism; moralism; autonomism; immoralism; merited response; imaginative resistance}

\begin{document}
\maketitle

\begin{abstract}
Can a work be morally flawed and aesthetically good? We defend a moderate moralism: a moral flaw in a work is an aesthetic flaw when, and only when, it prevents the response the work prescribes; otherwise it is a moral flaw only. The argument turns on the observation that many works aim at responses, such as pity, admiration or amusement, that are merited only if the work's moral perspective is acceptable. When a morally sensitive audience cannot give the prescribed response because the perspective is not acceptable, the work fails on its own terms, and that failure is aesthetic. We distinguish the position from Gaut's ethicism, which holds that every moral flaw that bears on prescribed responses is an aesthetic flaw, and from moderate autonomism, which denies that moral flaws are ever aesthetic flaws as such. We answer the immoralist objection that moral flaws can be aesthetic merits, and identify the class of moral flaws that remain aesthetically inert.
\end{abstract}

\section{The question}

\emph{Triumph of the Will} is admired for its cinematography and condemned for its glorification of Nazism. The question whether its moral character affects its value as a film, and not only its value all things considered, is the question of the relation between ethical and aesthetic evaluation. Three positions dominate the literature. \emph{Autonomism} holds that moral evaluation is irrelevant to aesthetic evaluation; in its moderate form, it allows that works may be morally assessed, but denies that a moral flaw is ever an aesthetic flaw as such \citep{AndersonDean1998}. \emph{Moralism} holds that moral flaws are, at least sometimes, aesthetic flaws. In its moderate form \citep{Carroll1996}, it holds that this is so in some cases; in the stronger form called ethicism \citep{Gaut1998}, it holds that it is so whenever the flaw is manifested in the attitudes a work prescribes. \emph{Immoralism} holds that moral flaws can be aesthetic merits, because a work that leads us to understand or imaginatively occupy a morally defective perspective can be the better for it \citep{Jacobson1997,Kieran2003}.

We defend a version of moderate moralism and give a criterion for the cases in which moral flaws are aesthetic flaws.

\section{Framework}

\begin{definition}\label{d:prescribed}
A response is \emph{prescribed} by a work if the work is designed so that an audience who understands it will have that response, and its success as the kind of work it is depends on their having it.
\end{definition}

Tragedies prescribe pity, comedies amusement, thrillers suspense, and works of political celebration admiration for what they celebrate. Prescription is fixed by the work's design, including its genre, not by the actual reactions of particular audiences.

\begin{definition}\label{d:merited}
A prescribed response is \emph{merited} if the features of the work's world to which it is directed give an audience reason to have it, given the perspective from which the work presents them.
\end{definition}

\begin{principle}[Uptake]\label{p:uptake}
A work fails as the kind of work it is, to the extent that a sensitive and informed audience cannot give the responses it prescribes.
\end{principle}

\ref{p:uptake} is the core of Carroll's argument \citep{Carroll1996}: works that prescribe responses are, in this respect, like jokes that do not amuse or thrillers that do not thrill. Failure to secure uptake is a failure of design, and failures of design are aesthetic failures. The qualification \enquote{sensitive and informed} excludes audiences who fail to respond through inattention or ignorance.

\begin{principle}[Moral sensitivity]\label{p:sensitive}
A sensitive audience does not give a response that would require it to endorse a perspective it recognizes as morally unacceptable.
\end{principle}

\ref{p:sensitive} is not a claim about what audiences ought to feel but about what they can: a morally sensitive person cannot be amused by a joke whose humour depends on contempt for the victims of an atrocity, any more than she can admire a celebration of the atrocity itself. This is related to the phenomenon of imaginative resistance \citep{Walton1994,Gendler2000}, discussed in the companion manuscript AES-08-b.

\section{The argument}

\begin{theorem}\label{t:main}
A moral flaw in a work is an aesthetic flaw if, and only if, it prevents a sensitive audience from giving a response the work prescribes.
\end{theorem}

\begin{proof}
Suppose a work's moral flaw consists in a perspective that a sensitive audience recognizes as unacceptable, and that a prescribed response requires endorsing that perspective. By \ref{p:sensitive}, the audience cannot give the response, and by \ref{p:uptake} the work fails to that extent as the kind of work it is. The flaw is therefore an aesthetic flaw. Conversely, suppose a moral flaw does not bear on any prescribed response: the work's success does not depend on the audience's endorsing the flawed perspective. Then the flaw does not prevent uptake, and by \ref{p:uptake} it does not make the work fail on its own terms; it is a moral flaw only.
\end{proof}

Theorem~\ref{t:main} differs from ethicism in its scope. Gaut holds that whenever a work manifests an ethically bad attitude in its prescribed responses, it is to that extent aesthetically defective. Our criterion is narrower: what matters is not whether the prescribed response is ethically bad, but whether it is blocked. A work may prescribe a response that is ethically questionable without blocking it, for instance because the flaw is subtle enough that sensitive audiences do not recognize it, or because the response does not require endorsing the flawed perspective but only entertaining it. Such a work is morally flawed and, on our view, not thereby aesthetically flawed, though it may become so for an audience that comes to recognize the flaw.

Two consequences follow. First, the aesthetic relevance of a moral flaw is relative to the audience's recognition of it, so that a work's aesthetic value can change as moral understanding changes. This is a familiar phenomenon: comedies whose humour depends on attitudes that were once widely shared have become unfunny, and this is an aesthetic loss for them. Second, moral flaws in works that do not prescribe responses dependent on their moral perspective, such as an abstract composition commissioned to glorify a regime, are not aesthetic flaws, since there is no response for them to block.

\section{Cases}

\emph{Celebration.} Works of celebration, such as \emph{Triumph of the Will}, are the clearest cases of aesthetic failure of moral origin, because the response they prescribe, admiration for what is celebrated, directly requires endorsing their perspective. The criterion classifies them uniformly: insofar as what they celebrate is recognized as evil, they fail as celebrations, whatever the merits of their means.

\emph{Comedy.} Humour is the response most sensitive to moral blockage, as \citet{Carroll1996} observed. A joke prescribes amusement, and amusement at a joke that depends on contempt for its target requires sharing the contempt, at least for the moment of the joke. When a sensitive audience recognizes the contempt as unwarranted, the joke fails, and the audience experiences the failure as a failure of humour, not only as a moral objection. This is why moral criticism of comedy so often takes the form \enquote{it isn't funny}.

\emph{Tragedy and sympathy.} A tragedy prescribes pity for its protagonist. If the work presents as sympathetic a character whose conduct a sensitive audience finds repugnant, and the work does not acknowledge the repugnance, pity is blocked and the work fails as a tragedy. AES-06 argues that the pleasure of tragedy lies in an understanding of loss achieved through pity and fear; the present criterion explains why a tragedy whose perspective is morally defective loses that pleasure as well, since there is no understanding to be had through a response one cannot give.

\emph{Incidental flaws.} Many works contain morally objectionable attitudes that play no role in the responses they prescribe: a stereotyped minor character in a novel whose concerns lie elsewhere, a casual remark in a letter. On our criterion these are moral flaws only. They may distract or irritate, and to that extent impede the work's other effects, but they do not make the work fail on its own terms. Ethicism counts them as aesthetic flaws if they are manifested in prescribed attitudes; we count them only when they block a prescribed response.

\section{Recent literature}

\citet{Gaut2007} developed ethicism into a full theory, with a cognitivist argument added to the merited-response argument: art can teach, and a work that teaches falsely about morality is to that extent aesthetically flawed. The cognitivist argument is independent of ours, and we do not rely on it. \citet{Giovannelli2007} mapped the positions in the debate by their answers to two questions, whether ethical and aesthetic value interact and whether the interaction is systematic, and placed moderate moralism among the views on which the interaction is real but not systematic. Our criterion makes the interaction systematic in one respect: it holds whenever a moral flaw blocks a prescribed response, and only then. It is in this sense a principled moderate moralism.

\section{Objections and replies}

\begin{objection}[Immoralism]
Some works are valuable because they lead us to occupy a morally defective perspective. Nabokov's \emph{Lolita} succeeds because its narrator's perspective is seductive. A moral flaw here is an aesthetic merit.
\end{objection}

\begin{reply}
The case does not involve a moral flaw in the work. \emph{Lolita} prescribes that we feel the seductiveness of Humbert's voice and recognize it as the voice of a predator; the work's perspective, as opposed to the narrator's, is not defective, and the prescribed response is not blocked for a sensitive reader. Immoralist cases generally have this structure: they are works that represent a defective perspective in order to illuminate it \citep{Kieran2003}. On our criterion, their moral value and their aesthetic value both lie in the illumination. A genuine immoralist case would be a work whose own perspective is defective and whose aesthetic success depends on the audience's endorsing it, and \ref{p:sensitive} implies that such a work cannot succeed with a sensitive audience.
\end{reply}

\begin{objection}[Autonomism]
When we judge \emph{Triumph of the Will} to be a flawed film, we are confusing two judgments: that it is a good film and that it is an evil one. The aesthetic judgment concerns its formal achievement, which is real \citep{AndersonDean1998}.
\end{objection}

\begin{reply}
The film prescribes admiration for what it celebrates. Its formal achievement is in the service of that response, and a sensitive audience cannot give it. On our criterion, it fails as a work of celebration, which is the kind of work it is, while succeeding in many of its formal means. This is not a confusion of two judgments but the recognition that a work can succeed in its means and fail in its end, which is an aesthetic judgment. Autonomism can describe the formal achievement but not the failure, which is why its verdict on the film seems incomplete.
\end{reply}

\begin{objection}[Audience relativity]
If the aesthetic relevance of a moral flaw depends on its recognition by the audience, the work's aesthetic value depends on who is watching. A film that succeeded with its original audience fails with a later one. This makes aesthetic value relative.
\end{objection}

\begin{reply}
The relevant audience is a sensitive and informed one, and sensitivity includes moral sensitivity. A flaw that a sensitive audience would recognize is a flaw in the work, even if the original audience did not recognize it, because they lacked the sensitivity in question. AES-03 argues that the verdicts of ideal critics converge; the ideal critic, in the sense relevant there, is also morally sensitive. The value of the work does not vary with the audience; what varies is how many actual audiences are in a position to see it.
\end{reply}

\section{What the result does not establish}

The criterion presupposes that the moral perspective of a work can be distinguished from that of its characters and narrators, and that its prescribed responses can be identified. Both are interpretive questions, and AES-09 argues that the author's intentions matter to them only when they succeed. The criterion does not tell us how to resolve hard interpretive cases; it tells us what follows once they are resolved. Nor does it say how much an aesthetic flaw of moral origin weighs against the work's other merits. A work may fail in one prescribed response and succeed in many others, and its overall value depends on their relative importance.

\section{Conclusion}

A moral flaw is an aesthetic flaw when it prevents a sensitive audience from responding as the work prescribes, because the work then fails on its own terms. When it does not, the flaw is moral only. The position is moderate in two directions: against autonomism, it holds that moral flaws can be aesthetic flaws; against ethicism, it holds that they are so only when they block uptake.

\begin{thebibliography}{99}\small
\bibitem[Anderson and Dean(1998)]{AndersonDean1998} Anderson, J.~C. and Dean, J.~T. (1998). Moderate autonomism. \emph{The British Journal of Aesthetics}, 38(2), 150--166.
\bibitem[Carroll(1996)]{Carroll1996} Carroll, N. (1996). Moderate moralism. \emph{The British Journal of Aesthetics}, 36(3), 223--238.
\bibitem[Gaut(1998)]{Gaut1998} Gaut, B. (1998). The ethical criticism of art. In J.~Levinson (Ed.), \emph{Aesthetics and Ethics: Essays at the Intersection} (pp.~182--203). Cambridge University Press.
\bibitem[Gaut(2007)]{Gaut2007} Gaut, B. (2007). \emph{Art, Emotion and Ethics}. Oxford University Press.
\bibitem[Gendler(2000)]{Gendler2000} Gendler, T.~S. (2000). The puzzle of imaginative resistance. \emph{The Journal of Philosophy}, 97(2), 55--81.
\bibitem[Giovannelli(2007)]{Giovannelli2007} Giovannelli, A. (2007). The ethical criticism of art: a new mapping of the territory. \emph{Philosophia}, 35(2), 117--127.
\bibitem[Jacobson(1997)]{Jacobson1997} Jacobson, D. (1997). In praise of immoral art. \emph{Philosophical Topics}, 25(1), 155--199.
\bibitem[Kieran(2003)]{Kieran2003} Kieran, M. (2003). Forbidden knowledge: the challenge of immoralism. In J.~L.~Bermúdez and S.~Gardner (Eds.), \emph{Art and Morality}. Routledge.
\bibitem[Walton(1994)]{Walton1994} Walton, K.~L. (1994). Morals in fiction and fictional morality. \emph{Proceedings of the Aristotelian Society, Supplementary Volume}, 68, 27--50.
\end{thebibliography}

\end{document}
